\documentclass[9pt,a4paper]{extarticle}

% ======================
% Codificação, idioma e layout
% ======================
\usepackage[T1]{fontenc}
\usepackage[utf8]{inputenc}
\usepackage[brazil]{babel}
\usepackage{geometry}
\geometry{margin=2cm}
\usepackage{parskip}
\usepackage{setspace}
\usepackage{microtype}
\usepackage{enumitem}

% ======================
% Matemática e símbolos
% ======================
\usepackage{amsmath, amssymb, amsthm, bm}

% Algoritmos (CLRS-style)
\usepackage{algpseudocode}

% Cores
\usepackage{xcolor}

% ======================
% Figuras e tabelas
% ======================
\usepackage{graphicx}
\usepackage{booktabs}
\usepackage{array}
\usepackage{tabularx}
\usepackage{float}
\usepackage{caption}
\usepackage{subcaption}

% ======================
% Hiperlinks
% ======================
\usepackage{hyperref}
\hypersetup{
  colorlinks=true,
  linkcolor=blue,
  urlcolor=blue
}

% ======================
% Início do documento
% ======================
\begin{document}

\begin{center}
  \Large \textbf{Prova 1 \\ Introdução à Análise de Algoritmos e Caracterização de Tempos de Execução} \\[4pt]
  \normalsize Disciplina: PCC104 - Projeto de Análise de Algoritmos
\end{center}

\vspace{1cm}

\begin{enumerate}
  \item Qual, ou quais os problemas de utilizar unidades de tempo, por exemplo, segundos para analisar o tempo de execuçãoo de algoritmos? Qual a a estratégia mais adequada para esta tarefa?
  \item Considere uma algoritmo que executa $2^n$ operções e outro que executa $100n^2$ operações. Qual é o aumento do número de operações para cada algortimo quando o tamanho da entrada cresce $4$ vezes?
  \item O que é uma \textit{invariante de laço} e como ela pode ser utilizada para provar a corretude de algoritmos?
  \item Considere a seguinte afirmação: \\
  \textit{A notação $O(.)$ indica a ordem de complexidade de pior caso de um algoritmo.}\\
  Esta afirmação é verdadeira ou falsa? Explique.
  \item Sejam $f(n)$ e $g(n)$ funções assintóticamente positivas, prove ou refute cada uma das seguintes conjecturas:
  \begin{enumerate}
    \item $f(n) = O(g(n))$ implica $g(n) = O(f(n))$
    \item $f(n) + g(n) = O(\max(g(n),f(n)))$
    \item $f(n) = \Theta(f(n/2))$
  \end{enumerate}
  \item Apresente a expresão completa do custo $T(n)$ do algoritmo abaixo:
    \begin{algorithmic}[1]
    \For{$i \gets 1$ \textbf{to} $A.\text{length} - 1$}
        \For{$j \gets A.\text{length}$ \textbf{downto} $i + 1$}
            \If{$A[j] < A[j-1]$}
                \State \textbf{exchange} $A[j]$ \textbf{with} $A[j-1]$
            \EndIf
        \EndFor
    \EndFor
  \end{algorithmic}
  \item Apresente a expresão completa do custo $T(n)$ do algoritmo abaixo:
  \begin{algorithmic}[1]
  \For{$j \gets 2$ \textbf{to} $A.\text{length}$}
      \State $key \gets A[j]$
      \State $i \gets j - 1$
      \While{$i > 0$ \textbf{and} $A[i] > key$}
          \State $A[i + 1] \gets A[i]$
          \State $i \gets i - 1$
      \EndWhile
      \State $A[i + 1] \gets key$
  \EndFor
  \end{algorithmic}
  \item Apresente a expresão completa do custo $T(n)$ do algoritmo abaixo e depois discuta a complexidade de execução deste algoritmo no melhor e no pior caso:
  \begin{algorithmic}[1]
  \Procedure{Sequential-Search}{$A, n, v$}
      \For{$i \gets 1$ \textbf{to} $n$}
          \If{$A[i] = v$}
              \State \Return $i$
          \EndIf
      \EndFor
      \State \Return \textsc{nil}
  \EndProcedure
  \end{algorithmic}
  \end{enumerate}


\end{document}

